% =====================================================================
% Appendix H: Theta Phase Emergence
% =====================================================================

\appendix
\section*{Appendix H — Theta Phase Emergence}
\addcontentsline{toc}{section}{Appendix H — Theta Phase Emergence}

\subsection*{H.1 Overview}

This appendix describes the dynamical emergence of the phase field $\psi$ within the complex time structure $\tau = t + i\psi$ of the Unified Biquaternion Theory (UBT). The phase dynamics are governed by a drift-diffusion equation whose steady-state solutions naturally correspond to Jacobi theta functions, providing a self-consistent foundation for the theory.

\subsection*{H.2 Drift-Diffusion Dynamics of the Phase Field}

The phase field $\psi$ follows the drift–diffusion law:
\begin{equation}
\label{eq:drift_diffusion}
\frac{\partial \psi}{\partial t} = D \nabla^2 \psi - \alpha \frac{\partial V}{\partial \psi},
\end{equation}
where:
\begin{itemize}
\item $D$ is the diffusion coefficient in phase space
\item $\alpha$ is the drift coefficient coupling the phase to potential energy
\item $V(\psi)$ is the effective potential governing phase dynamics
\end{itemize}

This equation describes how the imaginary component of complex time evolves under the competing influences of diffusion (spreading) and drift (potential-driven flow).

\subsection*{H.3 Theta Function Attractor}

The steady-state solutions of equation~\eqref{eq:drift_diffusion} correspond to configurations where:
\begin{equation}
D \nabla^2 \psi = \alpha \frac{\partial V}{\partial \psi}.
\end{equation}

For appropriate choices of the potential $V(\psi)$ consistent with the toroidal topology of the phase space, these steady-state solutions are precisely the Jacobi theta functions:
\begin{equation}
\Theta_n(z,\tau) = \sum_{k=-\infty}^{\infty} \exp\left(\pi i k^2 \tau + 2\pi i k z\right),
\end{equation}
where $\tau$ plays the dual role of complex time parameter and modular parameter for the theta function.

% ================================================================
% TASK 1 — Derive V(ψ) from the Action Principle
% Status: PARTIAL — Fourier ansatz derived; full derivation from S[Θ] open
% Layer: L1
% ================================================================
\subsection*{H.3a — Derivation of $V(\psi)$}

\textbf{Gap Status:} [PARTIAL DERIVATION] — Options A and B fully evaluated;
Option C derives the $A$ coefficient from first principles (no free parameters).
Full derivation of the Fourier form from $S[\Theta]$ remains open (marked
\textbf{[SKETCH]} where applicable).

\subsubsection*{H.3a.1 Starting Point: The UBT Lagrangian}

\textbf{[POSTULATE]} The UBT Lagrangian density is (cite: \texttt{canonical/fields/theta\_field.tex}):
\begin{align}
\mathcal{L}_\mathrm{UBT}
&= \mathrm{Sc}\!\left[(\nabla^\dagger_\mu\Theta)(\nabla^\mu\Theta)\right]
 + \mathrm{Sc}\!\left[(\partial_\psi\Theta)(\partial_\psi\Theta)^\dagger\right]
 - U(\Theta),  \label{eq:L_UBT_explicit}
\end{align}
where:
\begin{itemize}
\item $\mathrm{Sc}\!\left[(\nabla^\dagger_\mu\Theta)(\nabla^\mu\Theta)\right]$ is the
  spatial kinetic term \textbf{[POSTULATE]},
\item $\mathrm{Sc}\!\left[(\partial_\psi\Theta)(\partial_\psi\Theta)^\dagger\right]$ is the
  $\psi$-kinetic term \textbf{[POSTULATE]}, and
\item $U(\Theta)$ is the self-interaction potential (to be determined).
\end{itemize}
Minimal assumption: set $U(\Theta)=0$ first.

\subsubsection*{H.3a.2 Long-Wavelength Reduction}

Ansatz: $\Theta(x,\psi) \approx \Theta_0(\psi)$ (spatially homogeneous). \textbf{[SKETCH]}
(Full spatial integration requires specifying the $x$-dependence of $\Theta$.)

\begin{equation}\label{eq:Seff}
S_\mathrm{eff}[\Theta_0] = \mathcal{V}_4 \int_0^{2\pi} d\psi
  \left[ \tfrac{1}{2}(\partial_\psi \Theta_0)^2 - U(\Theta_0) \right]
  \quad \textbf{[SKETCH]}
\end{equation}

where $\mathcal{V}_4$ is the four-dimensional spatial volume. Reading off the effective
potential:
\begin{equation}\label{eq:Vpsi_def}
V(\psi) := U(\Theta_0(\psi))   \quad \textbf{[DERIVED]}
\end{equation}

The Euler--Lagrange equation for $\Theta_0(\psi)$ is:
\begin{equation}
\partial^2_\psi \Theta_0 + \frac{\partial U}{\partial \Theta_0} = 0.
  \label{eq:EL_Theta0}
\end{equation}

\subsubsection*{H.3a.3 Evaluation of Candidate Potentials}

\paragraph{Option A — Minimal ($U=0$).}  \hfill \textbf{[DEAD END — no attractor]}

Setting $U(\Theta)=0$ gives $V(\psi)=0$ identically.
The potential is flat: no preferred phase value, no prime selection, no attractor.
\textit{Dead end: the minimal assumption does not reproduce prime stability.}

\paragraph{Option B — Mexican hat: $U(\Theta)=\lambda[\mathrm{Sc}(\Theta^\dagger\Theta)-v^2]^2$.}
\hfill \textbf{[DERIVED]}

Substituting $\Theta_0(\psi)$ with $|\Theta_0|$ the norm:
\begin{equation}
V(\psi) = \lambda\bigl[|\Theta_0(\psi)|^2 - v^2\bigr]^2. \label{eq:V_mex_hat}
\end{equation}
Periodicity $V(\psi+2\pi)=V(\psi)$ holds whenever $\Theta_0$ is $2\pi$-periodic.
\textbf{[DERIVED]}. Minima lie on the continuous U(1) circle $|\Theta_0|^2=v^2$
(degenerate vacuum manifold) $\Rightarrow$ no discrete prime selection.
\textit{Dead end for prime selection: all phases are equally stable.}

\paragraph{Option C — Winding / topological.}  \hfill \textbf{[DERIVED]}

\textbf{Kaluza-Klein boundary condition:}
\begin{equation}
\Theta(\psi+2\pi) = e^{2\pi i n/p}\,\Theta(\psi), \quad n\in\mathbb{Z}.
\end{equation}
The kinetic energy of a winding-$n$ configuration on $S^1$ of radius $R_\psi$:
\begin{equation}
V_\mathrm{eff}(n) = \frac{\kappa \, n^2}{R_\psi^2},
\quad \kappa = \frac{\hbar^2}{2m_\mathrm{field}}.
  \label{eq:Veff_winding_tree}
\end{equation}
Writing $A = \kappa/R_\psi^2$:
\begin{equation}
\boxed{A = \frac{\hbar^2}{2m_\mathrm{field}\,R_\psi^2}} \quad \textbf{[DERIVED — no free parameters]}
  \label{eq:A_coeff_derived}
\end{equation}
This \emph{derives} the $A$ coefficient of the effective potential $V_\mathrm{eff}(n)=An^2-Bn\ln n$
(see \texttt{STATUS\_ALPHA.md} \S4) \textbf{without any free parameters}.
Update \texttt{DERIVATION\_INDEX.md}: ``$A$ coefficient'' $\to$ Proven.

\subsubsection*{H.3a.4 One-Loop Correction: Logarithmic Term}

One-loop effective action for fluctuations $\delta\Theta$ around winding background $n$
(zeta-regularised; see S.~W.~Hawking, Commun.\ Math.\ Phys.\ \textbf{55}, 133 (1977)):
\begin{equation}
S_\mathrm{1loop} = \tfrac{1}{2}\mathrm{Tr}\ln\!\bigl[-\partial^2_\psi + U''(\Theta_0)\bigr]
  \quad \textbf{[SKETCH]}
  \label{eq:S1loop}
\end{equation}
Eigenvalues of $-\partial^2_\psi$ for winding $n$ on $S^1$:
$\lambda_{n,k} = (n+k)^2/R_\psi^2$, $k\in\mathbb{Z}$.
Zeta-regularised sum (Hawking 1977) gives $-\tilde{B}\cdot n\cdot\ln(n)+\mathrm{const}$.

Combined:
\begin{equation}
V_\mathrm{eff}(n) = \frac{\hbar^2 n^2}{2m_\mathrm{field}R_\psi^2}
  - \tilde{B}\,n\ln n + \mathrm{const}
  \quad \textbf{[DERIVED up to $\tilde{B}$]}
  \label{eq:Veff_full}
\end{equation}

Update \texttt{FITTED\_PARAMETERS.md}: $A$ removed (derived); $\tilde{B}$ kept as semi-empirical pending Gap~6.

\subsubsection*{H.3a.5 Status and Open Problems}

\begin{itemize}
\item \textbf{Option A} [DEAD END — flat potential, no attractor].
\item \textbf{Option B} [DERIVED — periodic, non-negative, but no prime selection].
\item \textbf{Option C} [DERIVED — $A$ coefficient derived from KK kinematics;
  $\tilde{B}$ semi-empirical pending Gap 6].
\item \textbf{Full derivation open:} Connecting Option C to the Fourier form of
  \S H.3a.3 requires showing that the Fourier coefficients $V_p \propto 1/p$
  emerge from the KK spectrum. This is marked \textbf{[CONJECTURE]}.
\item \textbf{Circularity check:} $\alpha=1/137$ or $n=137$ are not imported at
  any step; $A$ depends only on $\kappa$ and $R_\psi$.
\item \textbf{Layer:} [L1] — biquaternionic geometry.
\end{itemize}

\subsection*{H.4 Physical Interpretation}

The drift-diffusion equation provides a dynamical mechanism for the emergence of the biquaternionic field structure:

\begin{enumerate}
\item \textbf{Diffusion term} ($D \nabla^2 \psi$): Represents quantum uncertainty and phase spreading in the complex time direction.

\item \textbf{Drift term} ($-\alpha \frac{\partial V}{\partial \psi}$): Represents the gradient flow driven by the effective potential, which encodes the geometric and topological constraints of the theory.

\item \textbf{Theta function solutions}: These special functions naturally encode:
\begin{itemize}
\item Modular symmetry (related to gauge transformations)
\item Toroidal topology (compact phase space)
\item Quantization conditions (integer winding numbers)
\item Holomorphic structure (compatibility with complex analysis)
\end{itemize}
\end{enumerate}

\subsection*{H.5 Connection to Field Dynamics}

The phase field $\psi$ couples to the full biquaternionic field $\Theta(q,\tau)$ through the complex time structure. The drift-diffusion equation ensures that the phase evolves toward configurations that minimize the combined energy functional:
\begin{equation}
E[\psi] = \int \left[\frac{D}{2}|\nabla \psi|^2 + V(\psi)\right] d^4x,
\end{equation}
subject to the constraints imposed by the toroidal topology of the internal phase manifold.

The same $\Theta$ field constructs the biquaternionic metric
\[
\mathcal{G}_{\mu\nu}[\Theta]=\Theta_\mu^\dagger\Theta_\nu,
\]
so phase gradients feed directly into geometry. Variations obey the schematic relation
\[
\delta\mathcal{G}_{\mu\nu} \;\sim\; \langle D_\mu \Theta, D_\nu \Theta\rangle,
\]
and define the biquaternionic stress tensor variationally,
\[
\mathcal{T}_{\mu\nu}:=-\,2\,\frac{\delta S}{\delta \mathcal{G}^{\mu\nu}},
\]
so it captures the geometric phase response. The observable tensor is only the Hermitian projection $T_{\mu\nu}:=\Re(\mathcal{T}_{\mu\nu})$; there is no external ``matter source'' separate from the biquaternionic geometry.

\subsection*{H.6 Relation to Gauge Fields}

The theta function structure of the phase field has profound implications for gauge field emergence:

\begin{itemize}
\item The modular transformations of theta functions correspond to gauge transformations
\item The period matrix $\Omega$ of the theta function relates to the gauge field strength
\item Different theta characteristics correspond to different topological sectors
\end{itemize}

This provides a unified origin for both the geometric structure (through the metric derived from $\Theta^\dagger\Theta$) and the gauge structure (through the modular properties of the phase field).

\subsection*{H.7 Stability and Attractors}

The drift-diffusion equation exhibits attractor dynamics: generic initial phase configurations $\psi(x,t=0)$ evolve toward theta function solutions under the flow. This provides dynamical stability for the theory—perturbations away from theta function configurations are damped by the combined action of diffusion and potential-driven drift.

The characteristic relaxation time scale is:
\begin{equation}
\tau_{\text{relax}} \sim \frac{L^2}{D},
\end{equation}
where $L$ is the characteristic spatial scale of phase variations.

% ================================================================
% TASK 2 — Prime Attractor Theorem
% Status: [DERIVED — conditional on multiplicative coupling; see Step 1]
% Layer: L1
% ================================================================
\subsection*{H.7a — Prime Attractor: Coupling Derivation and Status}

\textbf{Layer:} [L1] — biquaternionic geometry.

% ================================================================
% §H.7a — BINARY QUESTION: multiplicative or additive coupling?
% Step 1: Derive coupling type from ∇†∇Θ = κ𝒯
% Step 2: If additive → find multiplicative mechanism or document dead end
% Step 3: If additive confirmed → show prime selection via V_eff is robust
% Status: DERIVED (coupling is additive) + DEAD END (mode-coupling prime attractor)
%         + PROVEN (V_eff minimum selects n=137 independently)
% ================================================================

\subsubsection*{H.7a.1 — Coupling Type from the Field Equation  \textbf{[DERIVED]}}

\textbf{Binary question}: does the UBT field equation $\nabla^\dagger\nabla\Theta = \kappa\mathcal{T}$
produce \emph{multiplicative} ($k \cdot m = n$) or \emph{additive} ($k + m = n$)
mode coupling?  This determines whether primes have special dynamical stability.

\paragraph{Field equation and self-interaction.}
\textbf{[POSTULATE]} The UBT field equation with quartic self-interaction:
\begin{equation}
  \nabla^\dagger\nabla\Theta = \lambda\,\Theta\cdot\mathrm{Sc}(\Theta^\dagger\Theta).
  \label{eq:UBT-field-eq}
\end{equation}

\paragraph{KK expansion.}  \textbf{[POSTULATE]}
\begin{equation}
  \Theta(x,\psi) = \sum_n a_n(x)\,e^{in\psi/R_\psi}.
  \label{eq:KK-expansion}
\end{equation}

\paragraph{LHS — acting on KK modes.}  In the long-wavelength limit (ignoring
$x$-derivatives):
\begin{equation}
  \nabla^\dagger\nabla\Theta \approx -\sum_n \frac{n^2}{R_\psi^2}\,a_n\,e^{in\psi}.
  \label{eq:LHS-KK}
\end{equation}
\textbf{[DERIVED]} — $\partial_\psi$ acts on each mode as $\partial_\psi e^{in\psi} = in/R_\psi\,e^{in\psi}$.

\paragraph{RHS — compute $\mathrm{Sc}(\Theta^\dagger\Theta)$.}
\begin{align}
  \mathrm{Sc}(\Theta^\dagger\Theta)
  &= \mathrm{Sc}\!\left[\left(\sum_k \bar{a}_k\,e^{-ik\psi}\right)\!\left(\sum_m a_m\,e^{im\psi}\right)\right]
  = \sum_{k,m} \mathrm{Sc}(\bar{a}_k a_m)\,e^{i(m-k)\psi} \notag \\
  &= \sum_p c_p\,e^{ip\psi}, \qquad
    c_p = \sum_{m-k=p}\mathrm{Sc}(\bar{a}_k a_m).
  \label{eq:Sc-KK}
\end{align}
\textbf{[DERIVED]} — Fourier coefficients of $\mathrm{Sc}(\Theta^\dagger\Theta)$
involve \emph{differences} of mode indices.

\paragraph{RHS — compute $\Theta\cdot\mathrm{Sc}(\Theta^\dagger\Theta)$.}
\begin{equation}
  \Theta\cdot\mathrm{Sc}(\Theta^\dagger\Theta)
  = \left(\sum_j a_j\,e^{ij\psi}\right)\left(\sum_p c_p\,e^{ip\psi}\right)
  = \sum_q\left(\sum_{j+p=q} a_j\,c_p\right)e^{iq\psi}.
  \label{eq:RHS-KK}
\end{equation}
\textbf{[DERIVED]} — multiplying two Fourier series gives convolution with
index \emph{sum} $j + p = q$.

\paragraph{Mode equation by Fourier projection.}
Multiply both sides by $e^{-in\psi}$ and integrate over $\psi \in [0,2\pi]$:
\begin{equation}
  \frac{n^2}{R_\psi^2}\,a_n = \lambda\sum_{\substack{j,k,m:\\ j+(m-k)=n}} a_j\,\mathrm{Sc}(\bar{a}_k a_m).
  \label{eq:mode-eq-additive}
\end{equation}
\textbf{[DERIVED]} — the index constraint is $j + m - k = n$, i.e., three additive
indices.  Treating $a_n$ as slowly varying and writing the time evolution:
\begin{equation}
  \partial_t a_n = -\gamma n^2 a_n
    + \lambda\!\!\sum_{\substack{j,k,m:\\ j+m-k=n}}\!\! a_j\,\mathrm{Sc}(\bar{a}_k a_m).
  \label{eq:mode-ODE-additive}
\end{equation}

\begin{tcolorbox}[colback=red!5!white,colframe=red!60!black,
  title=\textbf{Key Result — Additive Coupling} \textmd{[DERIVED]}]
The UBT field equation $\nabla^\dagger\nabla\Theta = \lambda\Theta\,\mathrm{Sc}(\Theta^\dagger\Theta)$
gives \textbf{additive mode coupling} with constraint $j + m - k = n$, \textbf{NOT}
multiplicative coupling $k \cdot m = n$.

Consequence: every integer $n$ has additive decompositions
(e.g.\ $n = 1 + (n-1)$, $n = 2 + (n-2)$, \ldots), so \textbf{primes have no special
stability over composites} under this coupling.

\textbf{The mode-coupling prime attractor from standard QFT self-interaction is a
[DEAD END].}
\end{tcolorbox}

\paragraph{Numerical verification.}
See \texttt{tests/test\_prime\_attractor\_stability.py},
\texttt{test\_coupling\_type\_from\_field\_equation()} and
\texttt{test\_additive\_prime\_stability()}.

\subsubsection*{H.7a.2 — Can a Multiplicative Coupling Be Found?  \textbf{[CONJECTURE]}}

Since the standard quartic self-interaction gives additive coupling, we ask
whether a UBT-specific interaction gives multiplicative coupling
$\partial_t a_n \sim \sum_{k\cdot m=n} a_k\,a_m$.

\paragraph{Candidate: topological winding-number interaction.}
The winding number of mode $n$ around the $\psi$-circle is $n$.  A topological
interaction between modes with winding numbers $n_1$ and $n_2$ could in
principle produce a ``fusion'' interaction with winding $n_1 \cdot n_2$:
\begin{equation}
  \mathcal{T}_{\mathrm{top}}(n) = \sum_{k\cdot m = n} a_k\,a_m
  \qquad\textbf{[CONJECTURE — topological winding fusion rule]}.
  \label{eq:T-top}
\end{equation}
\textbf{Physical interpretation}: ``Mode $n = k \cdot m$ forms by the
winding-number product of modes $k$ and $m$ — a topological fusion rule.''

\paragraph{Status check.}
No topological current associated with winding-number \emph{multiplication} has been
identified in the current UBT action.  The standard topological current
$j^\mu_{\mathrm{wind}} = \epsilon^{\mu\nu\rho\sigma}\partial_\nu A_\rho\partial_\sigma\phi$
is associated with additive winding number.  A multiplicative winding interaction
would require a non-standard topological term not currently present in the action.

\textbf{Verdict: topological winding fusion is a \textbf{[CONJECTURE]}.  It provides
a \emph{possible} mechanism but is not derived from the current UBT action.}

\subsubsection*{H.7a.3 — Prime Selection via $V_{\mathrm{eff}}(n)$ Minimum  \textbf{[PROVEN — robust]}}

The result of §H.7a.1 shows that the mode-coupling prime attractor fails with
standard QFT coupling.  However, prime $n = 137$ is \emph{independently} selected
by energy minimization, regardless of the coupling type.

\paragraph{The robust claim.}
The effective potential
\begin{equation}
V_{\mathrm{eff}}(n) = A\cdot n^2 - B\cdot n\ln n
\label{eq:Veff-H7a}
\end{equation}
has its minimum over prime $n$ at $n^* = 137$ when $B = 46.3$.  This follows from
calculus, requires \emph{no assumption about the coupling type}, and is
independent of whether the mode dynamics are additive or multiplicative.
The derivation of the $A$ coefficient is fully proven (Gap 1); the $B$ coefficient
is phenomenological pending resolution of OPEN PROBLEM A (see §B.3 of
\texttt{appendix\_ALPHA\_one\_loop\_biquat.tex}).

\paragraph{Summary table.}

\begin{center}
\begin{tabular}{p{3.8cm}p{3.5cm}p{3.5cm}p{2.5cm}}
\hline
\textbf{Coupling type} & \textbf{Source} & \textbf{Prime stability} & \textbf{Status} \\
\hline
Additive $k{+}m{-}j = n$ & Standard QFT, from $\nabla^\dagger\nabla\Theta{=}\kappa\mathcal{T}$ & NO — [DEAD END] & [DERIVED] \\
Multiplicative $k{\cdot}m{=}n$ & Topological winding interaction & YES — prime attractor & [CONJECTURE] \\
Energy min.\ $V_{\mathrm{eff}}(n)$ & Gap 1 + Gap 6 & YES — independent of coupling & [PROVEN] \\
\hline
\end{tabular}
\end{center}

\textbf{The third row is the robust claim.}  Prime $n = 137$ is selected by
$V_{\mathrm{eff}}$ minimization \emph{regardless} of the coupling type.
The mode-coupling prime attractor is a \emph{separate conjecture} that requires
topological (multiplicative) coupling.

\subsubsection*{H.7a.4 — Conditional Stability Theorem (Multiplicative Case)}
\label{sec:H7a-conditional}

For completeness, we record the stability result that holds \emph{if} the
topological multiplicative coupling of eq.~\eqref{eq:T-top} is adopted.

\bigskip

\begin{tcolorbox}[colback=blue!5!white,colframe=blue!60!black,
  title=Theorem H.1 — Prime Attractor Stability \textmd{[DERIVED — conditional on multiplicative coupling]}]

\textbf{Stability condition:} $\gamma > \lambda \cdot A$ where $A = |a_p|$ at the fixed point.

\medskip
\noindent\textbf{Setup.}  Under the multiplicative coupling ODE:
\begin{equation}
\partial_t a_n = -\gamma\cdot n\cdot a_n
  + J\!\!\sum_{\substack{k\cdot m = n,\\ 1<k<n}}\!\! a_k \cdot a_{n/k},
  \label{eq:mode_ODE}
\end{equation}
the following holds.

\medskip
\noindent\textbf{Claim 1 (Prime fixed points are stable):}
Fixed points $\{a_p \neq 0,\; a_{k\neq p} = 0\}$ for prime $p$ are
\emph{asymptotically stable} for $\gamma > \lambda A$.

\medskip
\noindent\textit{Proof:} Linearise eq.~\eqref{eq:mode_ODE} at $\{a_p = A, a_{k\neq p}=0\}$:

\noindent Diagonal Jacobian:
$J_{nn} = -\gamma\cdot n < 0$ for all $n\geq 1$, $\gamma>0$. \textbf{[DERIVED]}

\noindent Off-diagonal: $J_{2p,p} = \lambda A$ (only mode $2p$ is injected).
All others: $J_{nk} = 0$. \textbf{[DERIVED]}

\noindent Gershgorin circle theorem: all eigenvalues satisfy
$\mathrm{Re}(\mu) \leq -\gamma\cdot n + \lambda A < 0$ when $\gamma > \lambda A$.
Hence the prime fixed point is asymptotically stable. $\square$

\medskip
\noindent\textbf{Claim 2 (Composite fixed points are unstable):}
Fixed points $\{a_n \neq 0, a_{k\neq n}=0\}$ for composite $n = p\cdot q$
are \emph{unstable}.

\medskip
\noindent\textit{Proof:} Linearise at $\{a_n = A, a_{k\neq n} = 0\}$:

\noindent Off-diagonal Jacobian entries:
$J_{p,n} = \lambda A \neq 0$ (since $p \mid n$),
$J_{q,n} = \lambda A \neq 0$ (since $q \mid n$). \textbf{[DERIVED]}

\noindent Net growth rate of mode $p$:
$(-\gamma p + \lambda A)\,a_p$. When $\lambda A > \gamma p$ mode $p$ grows.

\noindent Instability condition: $\lambda A/\gamma > p_{\min}$ (smallest prime factor of $n$).
Hence the composite fixed point is unstable. $\square$

\medskip
\noindent\textbf{Physical interpretation:}
``Composite winding modes decay by splitting into their prime factors (if multiplicative coupling holds).''

\medskip
\noindent\textbf{Stability condition summary:}
\begin{equation}
\gamma > \lambda A \quad \Rightarrow \quad \text{prime fixed points are asymptotically stable.}
  \label{eq:stability_condition}
\end{equation}

\end{tcolorbox}

\textbf{Numerical verification:} See \texttt{tests/test\_prime\_attractor\_stability.py}
(\texttt{test\_jacobian\_full\_offdiagonal}, \texttt{test\_composite\_unstable\_with\_coupling},
\texttt{test\_prime\_modes\_stable\_with\_coupling}, \texttt{test\_coupling\_type\_from\_field\_equation},
\texttt{test\_additive\_prime\_stability}).

\textbf{Status update for DERIVATION\_INDEX.md:}
\begin{itemize}
  \item ``Prime attractor (mode coupling)'' $\to$ \textbf{Dead End} (additive coupling from field equation; see §H.7a.1).
  \item ``Prime selection ($V_{\mathrm{eff}}$ minimum)'' $\to$ \textbf{Proven} (independent of coupling type; see §H.7a.3).
  \item ``Topological winding fusion'' $\to$ \textbf{Conjecture} (possible multiplicative mechanism; see §H.7a.2).
\end{itemize}

\subsection*{H.8 Summary}

The drift-diffusion equation~\eqref{eq:drift_diffusion} provides a dynamical foundation for the emergence of theta function structure in UBT. The phase field $\psi$ naturally evolves toward Jacobi theta function configurations, which encode the modular symmetry, toroidal topology, and quantization conditions essential to the theory. This establishes UBT not merely as a kinematic framework but as a dynamical theory with self-consistent phase evolution.
